La spécialité est, de tous les chapitres, celui qui ressemble le plus à ce qu'une bibliothèque
de mathématiques contient déjà. Division euclidienne, congruences, algorithme d'Euclide,
Bézout, Gauss, petit théorème de Fermat : chacun de ces résultats existe, souvent sous un
autre nom et dans une généralité plus grande. L'exercice n'est donc plus de démontrer mais de
retrouver la formulation exacte et de vérifier qu'elle dit bien la même chose que l'énoncé
français — exercice différent, et pas toujours plus facile.

La seconde moitié, matrices et graphes, est plus autonome. On y démontre l'inverse d'une
matrice $2 \times 2$, la résolution d'un système par l'inverse, les puissances par
diagonalisation, et la convergence d'un graphe probabiliste à deux états vers son état stable
— ce dernier résultat étant admis au lycée alors qu'il se démontre en trois lignes : l'écart à
l'état stable est multiplié à chaque étape par $1 - a - b$, dont le module est strictement
inférieur à $1$.

Un seul énoncé n'est pas formalisé : la correction du chiffrement RSA. Ses trois ingrédients —
petit théorème de Fermat, restes chinois, arithmétique modulaire — sont disponibles ; c'est
leur assemblage, et le traitement des cas dégénérés, qui demanderait un travail à part.
